%! Parent Functions and Transformations — Unit 2, Lesson 2.2 — Student Packet Cover Sheet
\documentclass[10pt]{article}
\usepackage{atda-article}
\usepackage{atda-boxes}
\usepackage{ltablex}
\keepXColumns

\begin{document}

% Full-bleed royal banner
\begin{tikzpicture}[remember picture, overlay]
  \fill[royal] (current page.north west)
    rectangle ([yshift=-0.9in]current page.north east);
\end{tikzpicture}
\vspace{-0.8in}

\begin{center}
  {\color{white}\textbf{\LARGE Algebra, Trigonometry, and Data Analysis}}\\[4pt]
  {\color{mistmid}\large\textbf{Unit 2: Functions, Transformations, and Graph Analysis}}\\[6pt]
  {\color{white}\textbf{\large Lesson 2.2 \quad Parent Functions and Transformations}}
\end{center}

\vspace{0.22in}
\namedateperiod
\vspace{0.18in}

\begin{learningtargetbox}
\small
\begin{enumerate}[leftmargin=1.5em, itemsep=5pt]
  \item \ldots{} identify the \textbf{parent function} of the linear, quadratic, and exponential
        families from an equation, a graph, or a table of values.
  \item \ldots{} describe the transformation that turns a parent function into a given function
        --- a \textbf{translation}, a \textbf{reflection}, or a \textbf{dilation} --- whether I am
        handed the equation or the graph.
  \item \ldots{} explain why a number \emph{outside} the parentheses moves a graph the way the sign
        says, while a number \emph{inside} moves it the opposite way.
  \item \ldots{} say which of the \textbf{domain} and the \textbf{range} a transformation moved,
        and use that as a check on my own answer.
\end{enumerate}
\end{learningtargetbox}

\vspace{0.14in}

\begin{tocbox}
\small
\renewcommand{\arraystretch}{1.5}
\begin{tabularx}{\linewidth}{@{} c l X r @{}}
\rowcolor{mist}
\textbf{\#} & \textbf{Component} & \textbf{Description} & \textbf{Score} \\
\midrule
1 & Warm-Up        & Feeding a rule a different number, and where a minus sign sits
                                                                    & \blank{1.2cm} \\
2 & Guided Notes   & The three parent functions, then the four moves --- and the card that
                     collects them                                  & \blank{1.2cm} \\
3 & Group Activity & Sliding a parabola, flipping an exponential, and a conveyor belt that speeds
                     up                                             & \blank{1.2cm} \\
4 & Exit Ticket    & Name the parent, name the move, and say which of the domain and range moved
                                                                    & \blank{1.2cm} \\
5 & Homework       & Families from tables, two irrigation zones, and a look ahead
                                                                    & \blank{1.2cm} \\
\midrule
  & \multicolumn{2}{r}{\textbf{Total}} & \blank{1.2cm} \\
\end{tabularx}
\end{tocbox}

\vspace{0.14in}

\begin{remindbox}
\small
This lesson covers \texttt{AFDA.AF.1a} (identify graphs and equations of the parent functions) and
\texttt{AFDA.AF.1b} (describe the transformation from the parent, given the equation or the graph).
One question settles every item in the packet: \textbf{is the number outside the parentheses, where
it changes the \emph{outputs}, or inside, where it changes the \emph{inputs}?} Outside moves the
graph the way the sign says. Inside moves it the opposite way --- and that is not a trick, it is
what substitution does. Bring Lesson 2.1's domain and range with you; they are your check.
\end{remindbox}

\end{document}
